% =========================================================================
% SENDA3 — guía educativa 3 (repositorio LATTIMEX)
% SPDX-License-Identifier: CC-BY-4.0
% Compilación: pdflatex guia_senda3 && bibtex guia_senda3 && pdflatex guia_senda3 && pdflatex guia_senda3
% Cifras y datos de las figuras en datos3/ (generados desde los datos crudos de los contratos).
% Esta guía no reemplaza a guia_senda.tex ni a guia_senda2.tex, que se conservan sin cambios.
% Edición pública, versión 1.2.0 del repositorio (motor lattimex-engine-1.8.0).
% =========================================================================
\documentclass[11pt,a4paper]{article}

\usepackage[T1]{fontenc}
\usepackage{lmodern}
\usepackage[spanish,es-noshorthands,es-nodecimaldot,es-tabla]{babel}
\usepackage[margin=2.4cm]{geometry}
\usepackage{microtype}
\usepackage{amsmath,amssymb}
\usepackage{booktabs}
\usepackage{array}
\usepackage{graphicx}
\usepackage{float}
\usepackage[font=small,labelfont=bf]{caption}
\usepackage{enumitem}
\usepackage[dvipsnames]{xcolor}
\usepackage{tikz}
\usetikzlibrary{arrows.meta,positioning,calc,shapes.geometric}
\usepackage{pgfplots}
\pgfplotsset{compat=1.18}
\usepackage[most]{tcolorbox}
\usepackage{algorithm}
\usepackage{algpseudocode}
\usepackage{url}
\usepackage[colorlinks=true,linkcolor=MidnightBlue,citecolor=OliveGreen,urlcolor=MidnightBlue]{hyperref}

\floatname{algorithm}{Algoritmo}
\algrenewcommand\algorithmicrequire{\textbf{Entrada:}}
\algrenewcommand\algorithmicensure{\textbf{Salida:}}
\algrenewcommand\algorithmicwhile{\textbf{mientras}}
\algrenewcommand\algorithmicdo{\textbf{hacer}}
\algrenewcommand\algorithmicif{\textbf{si}}
\algrenewcommand\algorithmicthen{\textbf{entonces}}
\algrenewcommand\algorithmicelse{\textbf{si no}}
\algrenewcommand\algorithmicfor{\textbf{para}}
\algrenewcommand\algorithmicend{\textbf{fin}}
\algrenewcommand\algorithmicreturn{\textbf{devolver}}
\algrenewtext{EndWhile}{\textbf{fin mientras}}
\algrenewtext{EndIf}{\textbf{fin si}}
\algrenewtext{EndFor}{\textbf{fin para}}

% --- Cajas didácticas (las mismas de las guías 1 y 2) ----------------------
\definecolor{sendaink}{RGB}{19,34,31}
\definecolor{sendared}{RGB}{210,75,50}
\definecolor{sendagreen}{RGB}{11,131,96}
\newtcolorbox{ideaclave}[1][]{enhanced,breakable,colback=sendagreen!6,colframe=sendagreen!70!black,
  fonttitle=\bfseries,title={Idea clave#1},left=6pt,right=6pt,top=4pt,bottom=4pt}
\newtcolorbox{ejemplo}[1][]{enhanced,breakable,colback=MidnightBlue!4,colframe=MidnightBlue!70,
  fonttitle=\bfseries,title={Ejemplo resuelto#1},left=6pt,right=6pt,top=4pt,bottom=4pt}
\newtcolorbox{leccion}[1][]{enhanced,breakable,colback=sendared!6,colframe=sendared!85!black,
  fonttitle=\bfseries,title={Lección metodológica#1},left=6pt,right=6pt,top=4pt,bottom=4pt}
\newtcolorbox{reproducir}[1][]{enhanced,breakable,colback=gray!6,colframe=gray!60!black,
  fonttitle=\bfseries,title={Para reproducir#1},left=6pt,right=6pt,top=4pt,bottom=4pt}

\newcommand{\senda}{SENDA}
\newcommand{\sendados}{SENDA2}
\newcommand{\sendatres}{SENDA3}
\newcommand{\pp}{\,pp}
\newcommand{\tsl}[1]{\texttt{\small #1}}

% Generado por gen_datos3.py desde los datos crudos de los contratos. No editar a mano.
\newcommand{\TresX}{0.79}
\newcommand{\DosX}{2.03}
\newcommand{\HgsX}{1.78}
\newcommand{\FiloX}{0.78}
\newcommand{\NXComunes}{97}
\newcommand{\TresABE}{0.065}
\newcommand{\DosABE}{0.062}
\newcommand{\HgsABE}{0.005}
\newcommand{\NABEComunes}{61}
\newcommand{\TresXG}{0.78}
\newcommand{\DosXG}{2.91}
\newcommand{\HgsXG}{2.55}
\newcommand{\FiloXG}{0.79}
\newcommand{\NXGComunes}{29}
\newcommand{\DXTresDos}{$-$1.25\pp{} [$-$1.46, $-$1.04]}
\newcommand{\DXTresDosCuenta}{89/3/5}
\newcommand{\DXTresDosN}{97}
\newcommand{\DXTresHgs}{$-$0.99\pp{} [$-$1.23, $-$0.76]}
\newcommand{\DXTresHgsCuenta}{75/2/20}
\newcommand{\DXTresHgsN}{97}
\newcommand{\DXTresFilo}{$-$0.01\pp{} [$-$0.08, +0.07]}
\newcommand{\DXTresFiloCuenta}{40/3/57}
\newcommand{\DXTresFiloN}{100}
\newcommand{\DXDosHgs}{+0.26\pp{} [+0.16, +0.36]}
\newcommand{\DXChicasTresDos}{$-$0.24\pp{} [$-$0.42, $-$0.07]}
\newcommand{\DXChicasTresDosCuenta}{13/3/5}
\newcommand{\DXChicasTresHgs}{+0.10\pp{} [$-$0.05, +0.23]}
\newcommand{\DXChicasTresHgsMedia}{+0.10}
\newcommand{\DXChicasTresHgsCuenta}{5/2/14}
\newcommand{\DXChicasTresHgsPeor}{14}
\newcommand{\DXChicasTresHgsN}{21}
\newcommand{\DXChicasTresFilo}{+0.11\pp{} [$-$0.03, +0.22]}
\newcommand{\DXChicasTresFiloCuenta}{3/2/16}
\newcommand{\DXMediasTresDos}{$-$1.03\pp{} [$-$1.21, $-$0.86]}
\newcommand{\DXMediasTresDosCuenta}{47/0/0}
\newcommand{\DXMediasTresHgs}{$-$0.72\pp{} [$-$0.92, $-$0.54]}
\newcommand{\DXMediasTresHgsCuenta}{41/0/6}
\newcommand{\DXMediasTresFilo}{+0.01\pp{} [$-$0.10, +0.12]}
\newcommand{\DXMediasTresFiloCuenta}{19/1/27}
\newcommand{\DXGrandesTresDos}{$-$2.32\pp{} [$-$2.66, $-$2.00]}
\newcommand{\DXGrandesTresDosCuenta}{29/0/0}
\newcommand{\DXGrandesTresHgs}{$-$2.21\pp{} [$-$2.59, $-$1.84]}
\newcommand{\DXGrandesTresHgsCuenta}{29/0/0}
\newcommand{\DXGrandesTresFilo}{$-$0.11\pp{} [$-$0.25, +0.03]}
\newcommand{\DXGrandesTresFiloCuenta}{18/0/14}
\newcommand{\DXNuevasTresDos}{$-$1.79\pp{} [$-$2.19, $-$1.25]}
\newcommand{\DXNuevasTresDosCuenta}{6/0/0}
\newcommand{\DXNuevasTresHgs}{$-$1.42\pp{} [$-$1.97, $-$0.84]}
\newcommand{\DXNuevasTresHgsCuenta}{6/0/0}
\newcommand{\DXNuevasTresFilo}{$-$0.20\pp{} [$-$0.53, +0.08]}
\newcommand{\DXNuevasTresFiloCuenta}{5/0/1}
\newcommand{\DXSinDevTresDos}{$-$1.20\pp{} [$-$1.42, $-$1.00]}
\newcommand{\DXSinDevTresHgs}{$-$0.93\pp{} [$-$1.17, $-$0.70]}
\newcommand{\DXSinDevTresFilo}{$-$0.02\pp{} [$-$0.09, +0.06]}
\newcommand{\DXCortasTresDos}{$-$1.16\pp{} [$-$1.47, $-$0.85]}
\newcommand{\DXCortasTresHgs}{$-$0.59\pp{} [$-$0.88, $-$0.32]}
\newcommand{\DXCortasTresFilo}{$-$0.28\pp{} [$-$0.44, $-$0.13]}
\newcommand{\DXLargasTresDos}{$-$1.28\pp{} [$-$1.53, $-$1.03]}
\newcommand{\DXLargasTresHgs}{$-$1.13\pp{} [$-$1.42, $-$0.85]}
\newcommand{\DXLargasTresFilo}{+0.09\pp{} [+0.02, +0.16]}
\newcommand{\DABETresDos}{+0.003\pp{} [$-$0.003, +0.009]}
\newcommand{\DABETresDosCuenta}{8/44/9}
\newcommand{\DABETresHgs}{+0.060\pp{} [+0.029, +0.095]}
\newcommand{\DABETresHgsCuenta}{0/41/20}
\newcommand{\DXGTresDos}{$-$2.13\pp{} [$-$2.47, $-$1.78]}
\newcommand{\DXGTresDosCuenta}{29/0/0}
\newcommand{\DTresIVDosX}{$-$1.90\pp{} [$-$2.21, $-$1.59]}
\newcommand{\DTresIVDosXN}{29}
\newcommand{\SentinelaDos}{0.999}
\newcommand{\TresCorridas}{1\,206}
\newcommand{\EnsayosContrato}{4\,214}
\newcommand{\DosFallasX}{15}
\newcommand{\DosFallasXG}{8}
\newcommand{\HgsFallasX}{11}
\newcommand{\HgsFallasXG}{2}
\newcommand{\IterTres}{630\,000}
\newcommand{\IterDos}{1\,140}
\newcommand{\HashContrato}{c91014f31b4f}
\newcommand{\HashRaw}{44bfb4763249}
\newcommand{\HashNucleo}{07fd3abe6d20}
\newcommand{\CurTresIV}{0.95}
\newcommand{\CurHgsIV}{3.21}
\newcommand{\CurFiloIV}{1.09}
\newcommand{\DCurHgsIV}{$-$2.25\pp{} [$-$2.62, $-$1.89]}
\newcommand{\DCurHgsIVMedia}{$-$2.25}
\newcommand{\DCurHgsIVCuenta}{30/0/0}
\newcommand{\DCurFiloIV}{$-$0.16\pp{} [$-$0.31, $-$0.01]}
\newcommand{\DCurFiloIVCuenta}{17/1/14}
\newcommand{\DCurFiloCortasIVMedia}{$-$0.46}
\newcommand{\CurTresXXX}{0.57}
\newcommand{\CurHgsXXX}{1.80}
\newcommand{\CurFiloXXX}{0.64}
\newcommand{\DCurHgsXXX}{$-$1.23\pp{} [$-$1.52, $-$0.96]}
\newcommand{\DCurHgsXXXMedia}{$-$1.23}
\newcommand{\DCurHgsXXXCuenta}{31/0/1}
\newcommand{\DCurFiloXXX}{$-$0.07\pp{} [$-$0.16, +0.01]}
\newcommand{\DCurFiloXXXCuenta}{15/0/17}
\newcommand{\DCurFiloCortasXXXMedia}{$-$0.26}
\newcommand{\CurTresCXX}{0.46}
\newcommand{\CurHgsCXX}{1.09}
\newcommand{\CurFiloCXX}{0.45}
\newcommand{\DCurHgsCXX}{$-$0.63\pp{} [$-$0.84, $-$0.43]}
\newcommand{\DCurHgsCXXMedia}{$-$0.63}
\newcommand{\DCurHgsCXXCuenta}{28/1/3}
\newcommand{\DCurFiloCXX}{+0.01\pp{} [$-$0.08, +0.09]}
\newcommand{\DCurFiloCXXCuenta}{12/0/20}
\newcommand{\DCurFiloCortasCXXMedia}{$-$0.15}
\newcommand{\DCurFiloLargasCXXCuenta}{5/0/18}
\newcommand{\DCruceIVHgsXXX}{$-$0.81\pp{} [$-$1.08, $-$0.55]}
\newcommand{\DCruceXXXHgsCXX}{$-$0.52\pp{} [$-$0.74, $-$0.32]}
\newcommand{\SentinelaTres}{0.998}
\newcommand{\HashCurva}{530e4ec70eb5}
\newcommand{\HashRawCurva}{8f59b9a19563}


\title{\textbf{SENDA3: rutas de mil puntos en segundos}\\[0.4em]
\large Tercera guía educativa del motor LATTIMEX: por qué un bucle nuevo, cómo se midió\\
contra HGS-CVRP y FILO2 con presupuestos de hasta 60 segundos, y dónde no gana}
\author{Erick Martínez\\
\small LATTIMEX · investigación independiente · Querétaro, México\\
\small \texttt{contacto@lattimex.com} · \texttt{https://lattimex.com}}
\date{Noviembre de 2026 · versión 1.2.0}

\hypersetup{
  pdftitle={SENDA3: rutas de mil puntos en segundos. Guía educativa 3 del motor LATTIMEX},
  pdfauthor={Erick Martínez (LATTIMEX)},
  pdfsubject={Ruteo de vehículos con capacidad (CVRP) con presupuestos de tiempo cortos},
  pdfkeywords={CVRP, SISR, recocido simulado, búsqueda local, HGS-CVRP, FILO2, evaluación reproducible}
}

\makeatletter
\def\ps@licenciaguia{%
  \let\@oddhead\@empty
  \let\@evenhead\@empty
  \def\@oddfoot{\vtop{\offinterlineskip%
    \hbox to\textwidth{\hfil\thepage\hfil}%
    \vskip5pt
    \hbox to\textwidth{\hfil\scriptsize\textcopyright{} 2026 Erick Martínez (LATTIMEX).
      CC BY 4.0: \url{https://creativecommons.org/licenses/by/4.0/}\hfil}}}%
  \let\@evenfoot\@oddfoot
}
\makeatother

\begin{document}
\maketitle
\thispagestyle{licenciaguia}

\begin{abstract}
\noindent\textbf{En pocas palabras.} La segunda guía dejó a \sendados{} cerca de HGS-CVRP en instancias
de 200 a 430 nodos, pero lejos de FILO2 y cada vez más lejos al crecer la instancia. Esta tercera guía
explica por qué pasaba eso y presenta \sendatres, que lo corrige. \sendatres{} está pensado para
responder en segundos: todos los resultados de este documento usan presupuestos de hasta 60\,s en un
solo hilo, y es en ese régimen donde se compara.

\smallskip
\noindent \sendatres{} es un selector entre dos núcleos. Con menos de 125 nodos corre \sendados{} sin
cambios. Desde 125 nodos corre un bucle nuevo de destruir y reconstruir \emph{en sitio}: quita cadenas
de clientes vecinos, las reinserta con una regla voraz, acepta con recocido simulado y, cada mil
iteraciones, aplica la búsqueda local de \sendados{} sólo a las rutas que cambiaron. El costo de una
iteración deja de depender del tamaño de la instancia: en 4\,s hace unas \IterTres{} iteraciones donde
\sendados{} hacía \IterDos.

\smallskip
\noindent Con un contrato congelado antes de correr, semillas nuevas y verificación independiente, en
las 100 instancias X (101 a 1001 nodos) a 4\,s \sendatres{} queda a \TresX\,\% de la mejor solución
conocida, frente a \DosX\,\% de \sendados, \HgsX\,\% de HGS-CVRP y \FiloX\,\% de FILO2. Contra HGS la
diferencia pareada es \DXTresHgs; contra FILO2 es un empate (\DXTresFilo). En las 32 instancias de 500
nodos o más, la ventaja sobre HGS es de \DCurHgsIVMedia\pp{} a 4\,s y de \DCurHgsXXXMedia\pp{} a 30\,s.
\sendatres{} no gana en todo: con menos de 200 nodos HGS sigue delante por una décima de punto o menos,
y a FILO2 no lo supera.
\end{abstract}

\noindent\textbf{Palabras clave:} ruteo de vehículos; CVRP; destruir y reconstruir; SISR; recocido
simulado; presupuestos de tiempo cortos; evaluación reproducible; docencia en optimización.

\tableofcontents
\newpage

% =========================================================================
\section*{Cómo leer este documento}
\addcontentsline{toc}{section}{Cómo leer este documento}

Esta guía continúa a las dos anteriores (\tsl{paper/guia\_senda.pdf} y \tsl{paper/guia\_senda2.pdf}) y
no las reemplaza. La primera explica \senda{} pieza por pieza y las bases de una comparación justa; la
segunda presenta \sendados{} y el acomodo tridimensional. Aquí se supone ese conocimiento y sólo se
trata el ruteo clásico con capacidad (CVRP).

\begin{itemize}[leftmargin=1.5em]
\item La Sección~\ref{sec:porque} explica por qué \sendados{} no escalaba.
\item La Sección~\ref{sec:bucle} describe el bucle nuevo, paso a paso.
\item La Sección~\ref{sec:selector} explica por qué \sendatres{} es un selector y no un reemplazo.
\item Las Secciones~\ref{sec:medicion} y~\ref{sec:curva} presentan la medición contra HGS-CVRP y FILO2.
\item La Sección~\ref{sec:pruebalo} explica cómo probarlo.
\end{itemize}

Los recuadros tienen las mismas cuatro funciones que en las guías anteriores: \emph{idea clave},
\emph{ejemplo resuelto}, \emph{lección metodológica} y \emph{para reproducir}.

\paragraph{El régimen de este documento.} \sendatres{} se diseñó para un planificador que responde
mientras el usuario espera. Por eso todas las mediciones usan presupuestos de pared de 2 a 30 segundos
(un punto adicional a 120\,s se reporta como referencia), en un hilo y en una sola máquina. Las cifras
publicadas de HGS-CVRP y FILO2 corresponden a presupuestos de varios minutos; esta guía no afirma nada
sobre ese régimen.

% =========================================================================
\section{Introducción}
\label{sec:intro}

La segunda guía cerraba con una advertencia y una pista. La advertencia: FILO2 era ``claramente más
fuerte en instancias grandes con reloj corto''. La pista, en el trabajo futuro: localizar la búsqueda
en los vértices tocados, ``la diferencia de diseño más probable detrás de la brecha en X''.

La meta de \sendatres{} se fijó en una frase: \textbf{meter instancias más grandes y resolverlas en
menos tiempo sin perder calidad}. El tamaño objetivo fue de 1\,000 puntos. La respuesta corta tiene
tres partes:

\begin{enumerate}[leftmargin=1.5em]
\item \textbf{El problema no eran los parámetros.} En una instancia de 1\,001 nodos, \sendados{} hacía
unas 210 iteraciones en 4\,s. Ajustes que multiplicaron las iteraciones por 18 no movieron el costo.
\item \textbf{Un bucle en sitio.} El bucle nuevo no copia la solución ni recorre estructuras del
tamaño de la instancia en cada iteración. Hace cientos de miles de iteraciones baratas en lugar de
cientos de iteraciones caras.
\item \textbf{Un selector.} En instancias chicas el bucle nuevo pierde por poco contra \sendados. En
lugar de forzarlo, \sendatres{} usa cada núcleo donde es mejor.
\end{enumerate}

\begin{ideaclave}
Cuando un algoritmo no escala, conviene preguntar primero cuánto cuesta una iteración y de qué
depende ese costo, antes de ajustar qué hace cada iteración.
\end{ideaclave}

% =========================================================================
\section{Por qué \sendados{} no escalaba}
\label{sec:porque}

\sendados{} es un ALNS: en cada iteración copia la solución, destruye entre 10\,\% y 35\,\% de los
clientes, los reinserta con una reparación \emph{regret}, mejora el resultado con búsqueda local y
decide si lo acepta. Con 100 clientes eso es rápido. Con 1\,000, destruir el 10\,\% significa quitar y
reinsertar 100 clientes, y la búsqueda local recorre toda la solución.

En una instancia de 1\,001 nodos, \sendados{} hacía unas 210 iteraciones en 4\,s. Se probó lo evidente:
destruir menos clientes, limitar la búsqueda local a lo tocado, cambiar la temperatura. Las iteraciones
se multiplicaron por 18 y el costo quedó en el mismo rango. Cada iteración seguía pagando trabajo
proporcional al tamaño de la instancia (copiar la solución, reconstruir índices, ordenar), y quitar
clientes sueltos producía muchas veces la misma solución, porque cada uno volvía a su sitio.

De ahí salieron las dos ideas de \sendatres:

\begin{itemize}[leftmargin=1.5em]
\item \textbf{Costo:} si una iteración sólo toca las rutas que cambia, su costo no depende del
número de clientes.
\item \textbf{Vecindario:} quitar \emph{cadenas} de clientes contiguos en rutas vecinas produce
cambios que la reinserción no deshace, a diferencia de quitar clientes sueltos.
\end{itemize}

Las dos apuntan al mismo diseño: el de SISR, \emph{Slack Induction by String
Removals}~\cite{christiaens2020}, que la segunda guía ya había usado en Python para el ruteo con
revisión de carga.

% =========================================================================
\section{El bucle de \sendatres}
\label{sec:bucle}

\sendatres{} arranca igual que \sendados: construye una solución y la pule con búsqueda local. Después,
en lugar del ALNS, repite cuatro pasos sobre una sola solución que modifica en sitio
(Algoritmo~\ref{alg:s3}).

\begin{algorithm}[H]
\caption{Bucle de \sendatres{} (esquema)}
\label{alg:s3}
\begin{algorithmic}[1]
\Require solución inicial $S$ (de \sendados), vecinos granulares, tiempo $t$
\State $T_0 \gets$ longitud media de un arco de $S$;\quad $T_{\text{fin}} \gets 0.01\,T_0$
\While{no se haya usado el 90\,\% de $t$}
  \State $T \gets T_0\,(T_{\text{fin}}/T_0)^{\text{fracción de tiempo usada}}$
  \State \textbf{destruir:} elegir un cliente semilla; quitar una cadena de la ruta de la semilla y de
         las rutas de sus vecinos, guardando una copia sólo de esas rutas
  \State \textbf{reconstruir:} ordenar los clientes quitados; insertar cada uno en la posición más
         barata junto a sus vecinos donde quepa
  \If{algún cliente no cabe en ninguna ruta} restaurar las rutas guardadas y continuar \EndIf
  \State \textbf{aceptar} si $\Delta < -T\ln u$, con $u$ uniforme en $(0,1)$; si no, restaurar las rutas guardadas
  \State cada 1\,000 iteraciones: búsqueda local de \sendados{} sobre los clientes de las rutas que cambiaron
\EndWhile
\State pulir la mejor solución con la búsqueda local de \sendados{} durante el 10\,\% restante
\end{algorithmic}
\end{algorithm}

\subsection{Destruir: cadenas en rutas vecinas}

Se elige un cliente semilla al azar. De su ruta se quita una cadena de clientes consecutivos que lo
contiene, de longitud aleatoria entre 1 y 10. Luego se recorren los vecinos más cercanos de la
semilla: de la ruta de cada uno (si no se ha tocado ya en esta iteración) se quita otra cadena que
contiene a ese vecino. El número de cadenas se sortea de modo que en promedio se quiten unos 10
clientes por iteración, sin importar si la instancia tiene 200 o 1\,000.

Quitar cadenas contiguas en rutas cercanas abre huecos \emph{alineados}: quedan varios extremos de ruta
próximos entre sí, y la reinserción puede cruzar tramos de una ruta a otra. Quitar clientes sueltos
deja huecos aislados, y lo más barato suele ser volver a poner cada cliente donde estaba.

\subsection{Reconstruir: inserción voraz con parpadeo}

Los clientes quitados se ordenan con una de cuatro reglas elegida al azar: orden aleatorio (4 de cada
11 veces), demanda descendente (4 de 11), los más lejanos al depósito primero (2 de 11) o los más
cercanos primero (1 de 11). Cada cliente se inserta en la posición que menos aumenta la distancia,
buscando sólo junto a sus 25 vecinos más cercanos y sólo en rutas donde cabe. Si ninguna de esas
posiciones sirve, se revisan todas las rutas; si aun así no cabe y la flota lo permite, se abre una
ruta.

El \emph{parpadeo}: al comparar posiciones, cada una se ignora con probabilidad de 1\,\%. Así la
reconstrucción no es siempre la misma para el mismo hueco, sin dejar de ser casi voraz.

\subsection{Aceptar: recocido con temperatura en arcos}

El cambio se acepta si empeora la distancia menos que un umbral aleatorio que depende de la temperatura
$T$. La temperatura baja de forma exponencial con el tiempo usado, de $T_0$ a $0.01\,T_0$, y $T_0$ es
la longitud media de un arco de la solución inicial. El criterio de aceptación y el enfriamiento
exponencial son los de SISR~\cite{christiaens2020}; medir la temperatura en múltiplos del arco medio
de la solución es la escala que usan FILO y FILO2~\cite{accorsi2021,accorsi2023}.

\begin{ejemplo}[: qué significa una temperatura de un arco]
Una solución de 1\,000 clientes y 43 rutas con distancia 75\,000 tiene $1\,000 + 43$ arcos, así que el
arco medio mide unas 72 unidades y $T_0 = 72$. Al inicio, un cambio que alarga la solución en 72
unidades se acepta con probabilidad $e^{-1}$, cerca de 37\,\%; uno que la alarga en 144, con
$e^{-2}$, cerca de 14\,\%. Al final, con $T = 0.72$, sólo pasan cambios que empeoran menos de una o dos
unidades. Medir la temperatura en arcos hace que el mismo par de valores sirva para instancias con
distancias muy distintas.
\end{ejemplo}

Si el cambio se rechaza, se restauran sólo las rutas guardadas al destruir. No hay copia de la solución
completa en ningún punto del bucle; esa es la diferencia que hace barata la iteración.

\subsection{La búsqueda local periódica}

\begin{leccion}[: el bucle solo no bastaba]
La primera versión del bucle, sin búsqueda local, ganaba en la mayoría de las instancias de prueba y
perdía por mucho en unas pocas. La inserción voraz no hace intercambios dentro de una ruta ni entre dos
rutas, y hay instancias que los necesitan. Agregar, cada mil iteraciones, la búsqueda local de
\sendados{} restringida a los clientes de las rutas que cambiaron corrigió esos casos. Promediar habría
escondido el problema: se vio por instancia.
\end{leccion}

La búsqueda local periódica tiene un tope de 0.02\,s por llamada y capacidad dura. Su resultado
reemplaza a la solución en curso sólo si es mejor.

\subsection{Flota}

Con flota libre, el bucle no usa más rutas que la solución inicial. Si la solución inicial viene
sobrecargada (pasa con rutas muy cortas, de tres o cuatro clientes), puede abrir rutas hasta eliminar
la sobrecarga y fija el tope en ese momento. Con flota fija, el tope es el número de vehículos. La
inserción nunca aumenta la sobrecarga, un cambio que la reduce se acepta siempre y la mejor solución
sólo se actualiza cuando no hay sobrecarga.

\begin{ideaclave}
En 4\,s, sobre las X de 125 nodos o más, \sendatres{} hace una mediana de \IterTres{} iteraciones por
corrida; \sendados{} hace \IterDos. No son comparables una a una (la de \sendados{} mueve decenas de
clientes y la de \sendatres{} unos diez), pero el cambio de escala es el punto.
\end{ideaclave}

% =========================================================================
\section{Por qué un selector}
\label{sec:selector}

Antes de la medición de esta guía hubo una primera confirmación del bucle, con reglas escritas antes de
correr. En instancias X el bucle mejoró a \sendados{} con holgura. En las instancias chicas A/B/E con
flota fija a 2\,s quedó unas centésimas de punto por encima, y la regla pedía no perder.

\begin{leccion}[: una regla que falla también decide]
La regla de A/B/E se escribió antes de correr precisamente para no poder ignorarla después. La
diferencia es pequeña, pero consistente. El protocolo decía qué hacer si fallaba sólo esa regla:
conservar \sendados{} para las instancias chicas.
\end{leccion}

Faltaba saber dónde poner el corte, porque A/B/E termina en 101 nodos. Se midieron las instancias X de
101 a 195 nodos, con flota libre y con flota fija, y una regla fijada de antemano: el umbral es el
menor tamaño a partir del cual el bucle no pierde en ningún tramo.

El umbral quedó en \textbf{125 nodos} (clientes más depósito). \sendatres{} es entonces:

\begin{itemize}[leftmargin=1.5em]
\item con menos de 125 nodos, \sendados{} tal cual;
\item con 125 nodos o más, el bucle nuevo;
\item si \sendados{} termina sin una solución válida (le pasa con flotas muy apretadas), el bucle
nuevo corre con el tiempo que quede.
\end{itemize}

El umbral se eligió y se evaluó con las mismas 21 instancias. Es una decisión de ingeniería. Lo que sí está confirmado es el
comportamiento del selector completo, que se mide a continuación.

% =========================================================================
\section{La medición}
\label{sec:medicion}

\subsection{Diseño}

Antes de correr se congeló un contrato con los brazos, los bancos, las semillas, el análisis y las
reglas de decisión, junto con las huellas SHA-256 de binarios, fuentes, instancias y scripts. El
programa de la corrida se niega a empezar si algo cambió.

\paragraph{Brazos.} \sendatres{} (el selector, tal como se integra al repositorio); \sendados{} (el
publicado en la versión 1.1.0); HGS-CVRP~\cite{vidal2022hgs} y FILO2~\cite{accorsi2023}, compilados de
sus fuentes oficiales, los mismos binarios de la segunda guía.

\paragraph{Reloj.} Pared, desde antes de leer la instancia hasta después de verificar la solución.
Todos reciben el 80\,\% de la pared como presupuesto interno. Un hilo, un ensayo a la vez, todos los
brazos en la misma sesión y con el orden rotado por instancia y réplica.

\paragraph{Validez.} Una corrida cuenta si la solución, revisada fuera del solver, cubre a todos los
clientes una vez, respeta la capacidad y la flota, su costo recalculado coincide y llegó dentro de la
pared. Distancias redondeadas al entero (convención de CVRPLIB~\cite{cvrplib}).

\begin{table}[H]
\centering
\caption{Bancos del contrato. Semillas 21000, 21100, \ldots{} (nunca usadas). \EnsayosContrato{} ensayos.}
\label{tab:bancos}
\small
\begin{tabular}{@{}llrlr@{}}
\toprule
Banco & Instancias & Pared & Flota & Réplicas \\
\midrule
X a 4\,s & las 100 X de Uchoa et al.~\cite{uchoa2017} (101--1001 nodos) & 4\,s & libre & 5 \\
A/B/E a 2\,s & 61 instancias A, B y E (22--101 nodos) & 2\,s & fija ($\le k$) & 10 \\
X grandes a 10\,s & las 32 X de 500 nodos o más & 10\,s & libre & 3 \\
\bottomrule
\end{tabular}
\end{table}

X a 4\,s con flota libre y A/B/E a 2\,s con flota fija son el reloj y el formato de las dos guías
anteriores. FILO2 no entra en A/B/E porque no tiene opción de flota fija.

\paragraph{Análisis.} Por instancia, media de las réplicas; una instancia entra en una comparación
sólo si los dos solvers tienen todas sus corridas aceptadas. Se reporta la diferencia pareada de
brecha en puntos porcentuales (pp), su intervalo de confianza de 95\,\% por remuestreo~\cite{efron1979},
la prueba de Wilcoxon~\cite{wilcoxon1945} y el conteo de instancias donde \sendatres{} es
mejor/empata/peor (empate si la diferencia es de 0.005\pp{} o menos).

\paragraph{Reglas fijadas antes de correr.} (R1) en X, \sendatres{} mejora a \sendados{} en al menos
0.5\pp; (R2) en X, el límite superior del intervalo de \sendatres{} menos HGS no pasa de cero; (R3) en
A/B/E, \sendatres{} no pierde contra \sendados{} (límite superior $\le +0.01$\pp) y queda a no más de
$+0.10$\pp{} de HGS; (R4) en las grandes, la ventaja sobre \sendados{} se mantiene a 10\,s y \sendatres{}
con 4\,s no es peor que \sendados{} con 10\,s; (R5) todas las corridas de \sendatres{} son válidas.

\begin{leccion}[: lo que este banco ya había visto]
Las instancias no son nuevas para \sendatres. El bucle se calibró en 7 instancias X y ya se había
medido en 67 con otras semillas; sólo 6 (de 439 a 491 nodos) no se habían usado nunca. Esta medición
es una réplica con semillas nuevas y configuración congelada, no una prueba en instancias nunca
vistas. Se reportan aparte las 6 nuevas y el resultado sin las 7 de calibración.
\end{leccion}

\subsection{Resultados}

Las cinco reglas se cumplieron. La corrida fue una sola sesión de 3 horas 37 minutos, sin
interrupciones. El sentinela de velocidad (iteraciones de \sendados{} en las 41 X del banco frente a
la medición del día anterior con el mismo binario) dio \SentinelaDos.

\begin{table}[H]
\centering
\caption{Brecha media contra la mejor solución conocida, en las instancias donde todos los solvers del
banco tuvieron todas sus corridas aceptadas. Misma sesión, mismo reloj, un hilo.}
\label{tab:principal}
\begin{tabular}{lrrrr}
\toprule
Banco & \sendatres & \sendados & HGS-CVRP & FILO2 \\
\midrule
X, 4\,s (\NXComunes{} de 100) & \TresX\,\% & \DosX\,\% & \HgsX\,\% & \textbf{\FiloX\,\%} \\
X grandes, 10\,s (\NXGComunes{} de 32) & \textbf{\TresXG\,\%} & \DosXG\,\% & \HgsXG\,\% & \FiloXG\,\% \\
A/B/E, 2\,s (\NABEComunes) & \TresABE\,\% & \DosABE\,\% & \textbf{\HgsABE\,\%} & -- \\
\bottomrule
\end{tabular}
\end{table}

\paragraph{Comparaciones pareadas en X a 4\,s.}
\begin{itemize}[leftmargin=1.5em]
\item \textbf{Contra \sendados:} \DXTresDos; mejor/empate/peor \DXTresDosCuenta{} en \DXTresDosN{} instancias.
\item \textbf{Contra HGS:} \DXTresHgs; \DXTresHgsCuenta{} en \DXTresHgsN.
\item \textbf{Contra FILO2:} \DXTresFilo; \DXTresFiloCuenta{} en \DXTresFiloN. Es un empate: el
intervalo contiene al cero y FILO2 gana en más instancias.
\end{itemize}
Como referencia, \sendados{} contra HGS en la misma sesión: \DXDosHgs, el mismo orden que reportó la
segunda guía ($+0.33$\pp{} en las 41 X del banco).

\begin{figure}[H]
\centering
\begin{tikzpicture}
\begin{axis}[width=0.95\textwidth,height=7cm,xlabel={nodos},ylabel={brecha media (\%)},
  legend style={at={(0.5,1.03)},anchor=south,legend columns=4,draw=none},legend cell align=left,
  grid=major,grid style={gray!20},ymin=0,ymax=7,unbounded coords=discard]
\addplot[only marks,mark=o,mark size=1.4pt,gray] table[x=n,y=senda2]{datos3/x_senda3.dat};\addlegendentry{\sendados}
\addplot[only marks,mark=square*,mark size=1.4pt,MidnightBlue] table[x=n,y=hgs]{datos3/x_senda3.dat};\addlegendentry{HGS-CVRP}
\addplot[only marks,mark=triangle*,mark size=1.7pt,sendagreen] table[x=n,y=filo2]{datos3/x_senda3.dat};\addlegendentry{FILO2}
\addplot[only marks,mark=*,mark size=1.6pt,sendared] table[x=n,y=senda3]{datos3/x_senda3.dat};\addlegendentry{\sendatres}
\end{axis}
\end{tikzpicture}
\caption{Brecha media por instancia X a 4\,s (5 semillas por solver, misma sesión). \sendados{} y HGS
se alejan de la mejor solución conocida al crecer la instancia; \sendatres{} y FILO2 se mantienen cerca
de 1\,\%. Faltan los puntos de \sendados{} y HGS en tres instancias donde no siempre dieron solución.}
\label{fig:x}
\end{figure}

\paragraph{Por tamaño.} La ventaja no es uniforme (Tabla~\ref{tab:tamano}).

\begin{table}[H]
\centering
\caption{\sendatres{} menos cada solver en X a 4\,s, por tramo. Diferencia pareada en pp con IC95;
entre paréntesis, instancias donde \sendatres{} es mejor/empata/peor.}
\label{tab:tamano}
\footnotesize
\begin{tabular}{@{}p{2.8cm}p{4.0cm}p{4.0cm}p{4.0cm}@{}}
\toprule
Tramo & contra \sendados & contra HGS-CVRP & contra FILO2 \\
\midrule
101--199 nodos & \DXChicasTresDos{} (\DXChicasTresDosCuenta) & \DXChicasTresHgs{} (\DXChicasTresHgsCuenta) & \DXChicasTresFilo{} (\DXChicasTresFiloCuenta) \\
200--499 & \DXMediasTresDos{} (\DXMediasTresDosCuenta) & \DXMediasTresHgs{} (\DXMediasTresHgsCuenta) & \DXMediasTresFilo{} (\DXMediasTresFiloCuenta) \\
500--1001 & \DXGrandesTresDos{} (\DXGrandesTresDosCuenta) & \DXGrandesTresHgs{} (\DXGrandesTresHgsCuenta) & \DXGrandesTresFilo{} (\DXGrandesTresFiloCuenta) \\
\midrule
las 6 nunca usadas & \DXNuevasTresDos{} (\DXNuevasTresDosCuenta) & \DXNuevasTresHgs{} (\DXNuevasTresHgsCuenta) & \DXNuevasTresFilo{} (\DXNuevasTresFiloCuenta) \\
sin las 7 de calibración & \DXSinDevTresDos & \DXSinDevTresHgs & \DXSinDevTresFilo \\
\midrule
rutas cortas ($<6$ clientes) & \DXCortasTresDos & \DXCortasTresHgs & \DXCortasTresFilo \\
rutas largas & \DXLargasTresDos & \DXLargasTresHgs & \DXLargasTresFilo \\
\bottomrule
\end{tabular}
\end{table}

Tres lecturas:
\begin{itemize}[leftmargin=1.5em]
\item \textbf{Con menos de 200 nodos, HGS sigue delante:} \DXChicasTresHgsMedia\pp, y \sendatres{}
pierde en \DXChicasTresHgsPeor{} de \DXChicasTresHgsN{} instancias. La ventaja sobre HGS empieza en 200
nodos y crece con el tamaño.
\item \textbf{Contra FILO2 depende de la forma de las rutas:} con rutas cortas \sendatres{} es mejor;
con rutas largas FILO2 es mejor por una décima de punto.
\item \textbf{Quitar las instancias de calibración no cambia el resultado,} y en las seis que no se
habían usado nunca \sendatres{} gana a \sendados{} y a HGS en las seis.
\end{itemize}

\paragraph{A/B/E a 2\,s.} \sendatres{} menos \sendados: \DABETresDos{} (\DABETresDosCuenta). Las 61
instancias tienen menos de 125 nodos, así que \sendatres{} ejecuta \sendados; la diferencia es ruido de
reloj entre dos ejecuciones del mismo algoritmo. Contra HGS: \DABETresHgs{} (\DABETresHgsCuenta);
HGS no pierde en ninguna. Este banco comprueba que el selector no rompe lo que ya funcionaba; no mide
el bucle nuevo.

\paragraph{Más grande en menos tiempo.} En las 32 X de 500 nodos o más, a 10\,s, \sendatres{} menos
\sendados{} es \DXGTresDos{} (\DXGTresDosCuenta). Y \sendatres{} con 4\,s es mejor que \sendados{} con
10\,s en las \DTresIVDosXN{} instancias comparables: \DTresIVDosX.

\paragraph{Validez.} Las \TresCorridas{} corridas de \sendatres{} fueron válidas y llegaron dentro de la
pared, lo mismo que todas las de FILO2. \sendados{} no dio solución válida en \DosFallasX{} de 500
corridas en X a 4\,s y en \DosFallasXG{} de 96 a 10\,s; HGS, en \HgsFallasX{} y \HgsFallasXG. Todas esas
fallas ocurren en tres instancias con rutas muy cortas, de tres a seis clientes (X-n524-k153, X-n670-k130 y
X-n936-k151), que por eso quedan fuera de las comparaciones pareadas con \sendados{} y HGS. Excluirlas
favorece a esos dos solvers, no a \sendatres.

\begin{reproducir}
Contrato \tsl{\HashContrato\ldots}; datos crudos \tsl{raw\_trials.jsonl} \tsl{\HashRaw\ldots}; núcleo
medido \tsl{\HashNucleo\ldots} (\tsl{native/senda3\_core.cpp} es ese fuente con los parámetros
fijos; da las mismas rutas). Con el repositorio:
\tsl{python -m lattimex cvrp INSTANCIA.vrp -{}-time 4 -{}-router senda3} (X) o
\tsl{-{}-time 2 -{}-fleet fixed -{}-router senda3} (A/B/E). \textbf{Corra una instancia a la vez}: el reloj
es de pared, y varios procesos en paralelo reducen las iteraciones de cada uno.
\end{reproducir}

% =========================================================================
\section{La curva de tiempo}
\label{sec:curva}

Una ventaja a 4\,s podría ser sólo un arranque más rápido. Para saber hasta dónde se sostiene dentro
del régimen de uso, un segundo contrato midió las 32 X de 500 nodos o más a 4\,s, 30\,s y 120\,s, con
\sendatres, HGS y FILO2 en la misma sesión, 2 réplicas y semillas nuevas (31000 y 31100). El punto de
120\,s queda fuera del régimen de 60\,s y se reporta como referencia.

\begin{table}[H]
\centering
\caption{Curva de tiempo en las 32 X de 500 nodos o más (flota libre, un hilo). Brecha media en las
instancias comunes y diferencia pareada con IC95.}
\label{tab:curva}
\small
\begin{tabular}{@{}lrrrp{4.0cm}p{4.0cm}@{}}
\toprule
Pared & \sendatres & HGS & FILO2 & \sendatres{} $-$ HGS & \sendatres{} $-$ FILO2 \\
\midrule
4\,s & \CurTresIV\,\% & \CurHgsIV\,\% & \CurFiloIV\,\% & \DCurHgsIV{} (\DCurHgsIVCuenta) & \DCurFiloIV{} (\DCurFiloIVCuenta) \\
30\,s & \CurTresXXX\,\% & \CurHgsXXX\,\% & \CurFiloXXX\,\% & \DCurHgsXXX{} (\DCurHgsXXXCuenta) & \DCurFiloXXX{} (\DCurFiloXXXCuenta) \\
120\,s & \CurTresCXX\,\% & \CurHgsCXX\,\% & \CurFiloCXX\,\% & \DCurHgsCXX{} (\DCurHgsCXXCuenta) & \DCurFiloCXX{} (\DCurFiloCXXCuenta) \\
\bottomrule
\end{tabular}
\end{table}

\begin{figure}[H]
\centering
\begin{tikzpicture}
\begin{axis}[width=0.75\textwidth,height=6.2cm,xmode=log,log basis x=10,xtick={4,30,120},
  xticklabels={4,30,120},xlabel={pared (s)},ylabel={brecha media (\%)},ymin=0,
  legend style={at={(0.5,1.03)},anchor=south,legend columns=3,draw=none},legend cell align=left,
  grid=major,grid style={gray!20}]
\addplot[mark=square*,MidnightBlue,thick] table[x=t,y=hgs]{datos3/curva_tiempo.dat};\addlegendentry{HGS-CVRP}
\addplot[mark=triangle*,sendagreen,thick] table[x=t,y=filo2]{datos3/curva_tiempo.dat};\addlegendentry{FILO2}
\addplot[mark=*,sendared,thick] table[x=t,y=senda3]{datos3/curva_tiempo.dat};\addlegendentry{\sendatres}
\end{axis}
\end{tikzpicture}
\caption{Brecha media contra el tiempo de pared en las 32 X de 500 nodos o más. La distancia entre
\sendatres{} y HGS se reduce al dar más tiempo, sin cerrarse en el rango medido.}
\label{fig:curva}
\end{figure}

\begin{itemize}[leftmargin=1.5em]
\item \textbf{Contra HGS,} el intervalo queda por debajo de cero en los tres tiempos. La ventaja se
reduce a cerca de la mitad en cada salto: \DCurHgsIVMedia, \DCurHgsXXXMedia{} y \DCurHgsCXXMedia\pp.
\item \textbf{En tiempo equivalente,} \sendatres{} con 4\,s es mejor que HGS con 30\,s
(\DCruceIVHgsXXX) y \sendatres{} con 30\,s es mejor que HGS con 120\,s (\DCruceXXXHgsCXX).
\item \textbf{Contra FILO2} es un empate a 30\,s y a 120\,s. A 4\,s el intervalo queda apenas bajo
cero, pero el conteo es \DCurFiloIVCuenta{} y la medición de la Sección~\ref{sec:medicion}, con más
réplicas, dio un intervalo que contiene al cero: no debe leerse como una ventaja.
\item \textbf{FILO2 aprovecha mejor el tiempo adicional.} Con rutas largas gana en más instancias a
120\,s (\sendatres{} mejor/empate/peor \DCurFiloLargasCXXCuenta); con rutas cortas \sendatres{}
conserva la ventaja en los tres tiempos (\DCurFiloCortasIVMedia, \DCurFiloCortasXXXMedia{} y
\DCurFiloCortasCXXMedia\pp).
\end{itemize}

Las 192 corridas de \sendatres{} fueron válidas; el sentinela de velocidad dio \SentinelaTres.

\begin{ideaclave}
Con el mismo reloj y un hilo, \sendatres{} entrega mejores rutas que HGS-CVRP a 4\,s desde 200 nodos,
y a 10\,s y 30\,s en las instancias de 500 nodos o más; con FILO2 empata. La curva dice además en qué dirección cambia eso: la
ventaja sobre HGS se reduce al dar más tiempo.
\end{ideaclave}

\begin{reproducir}
Contrato de la curva \tsl{\HashCurva\ldots}; datos crudos \tsl{\HashRawCurva\ldots}. Mismos binarios e
instancias que el contrato principal.
\end{reproducir}

% =========================================================================
\section{Discusión}
\label{sec:discusion}

\paragraph{Dónde está el valor.} \sendatres{} no es un algoritmo nuevo: combina el vecindario de
SISR~\cite{christiaens2020}, recocido simulado~\cite{kirkpatrick1983} con la temperatura medida en arcos, como en
FILO~\cite{accorsi2021}, y la búsqueda local granular que
\senda{} ya tenía~\cite{toth2003granular}. Lo propio es la combinación, la medición y un diseño que
cabe en unas 200 líneas sobre el núcleo anterior. Para quien planea rutas de cientos de entregas y
quiere la respuesta en segundos, la diferencia es práctica: en instancias de 500 a 1\,000 nodos, la
brecha a 4\,s baja de unos 3.3\,\% a cerca de 1\,\%.

\paragraph{Dentro del motor.} Desde la versión 1.2.0 el motor de LATTIMEX usa \sendatres{} como núcleo
de ruteo. El motor no llama al núcleo como el banco: reparte el tiempo entre ruteo, acomodo
tridimensional y reoptimización, y usa flota exacta. Conserva la regla que vuelve a \senda{} cuando hay
menos de 6 entregas por unidad. Esta guía mide sólo el ruteo clásico; no reporta cifras del motor con
acomodo tridimensional.

\paragraph{Limitaciones.}
(i) Régimen: presupuestos de 2 a 30\,s (y un punto a 120\,s), un hilo. No se afirma nada sobre
presupuestos de minutos ni sobre varios hilos.
(ii) Con menos de 200 nodos HGS-CVRP sigue siendo mejor, y en A/B/E no pierde en ninguna instancia.
(iii) A FILO2 no se le supera; con rutas largas y más tiempo, FILO2 es mejor.
(iv) Las instancias ya se habían usado durante el desarrollo, salvo seis.
(v) El umbral de 125 nodos se eligió con 21 instancias y no tiene confirmación propia.
(vi) Todas las instancias son de la familia X o de A/B/E; no hay otra familia entre 125 y 200 nodos.
(vii) No se reporta medición con flota fija en instancias X.
(viii) Todas las mediciones son de una sola máquina.

\section{Conclusiones y trabajo futuro}

\sendados{} no escalaba porque el costo de cada iteración crecía con la instancia. \sendatres{} lo
resuelve con un bucle en sitio de cadenas vecinas, una búsqueda local periódica y un selector que
conserva a \sendados{} donde sigue siendo mejor. Con contrato congelado y semillas nuevas, en las 100
instancias X a 4\,s baja la brecha de \DosX\,\% a \TresX\,\%, queda por debajo de HGS-CVRP con el mismo
reloj desde 200 nodos y empata con FILO2.

\paragraph{Trabajo futuro.} Cerrar la diferencia con HGS-CVRP por debajo de 200 nodos, y repetir la
medición en una segunda máquina y dentro del motor con acomodo tridimensional.

% =========================================================================
\section{Pruébelo en su computadora}
\label{sec:pruebalo}

\begin{verbatim}
git clone https://github.com/Erick-MF28/lattimex.git
cd lattimex
pip install -r requirements.txt  # Windows sin compilador: pip install ziglang
python -m lattimex build         # compila los núcleos y los acomodadores
python -m lattimex probe         # huella determinista
\end{verbatim}

\paragraph{\sendatres{} en CVRPLIB.}
\begin{verbatim}
python -m lattimex cvrp X-n1001-k43.vrp --time 4 --router senda3 --bks 72355
python -m lattimex cvrp X-n1001-k43.vrp --time 4 --router senda2 --bks 72355
python -m lattimex cvrp A-n32-k5.vrp --time 2 --fleet fixed --router senda3
\end{verbatim}
La primera y la segunda línea comparan los dos núcleos en su máquina. En la tercera, por tener menos
de 125 nodos, \sendatres{} ejecuta \sendados{} y lo indica en el campo \tsl{core} de la respuesta.

\begin{reproducir}[: artefactos]
\small
\begin{tabular}{@{}p{5.2cm}p{9.3cm}@{}}
Contrato principal & \tsl{contract.json} \tsl{\HashContrato\ldots}; \tsl{raw\_trials.jsonl} \tsl{\HashRaw\ldots} \\
Curva de tiempo & \tsl{contract.json} \tsl{\HashCurva\ldots}; \tsl{raw\_trials.jsonl} \tsl{\HashRawCurva\ldots} \\
Núcleo medido & \tsl{\HashNucleo\ldots}; \tsl{native/senda3\_core.cpp} es ese fuente con los parámetros fijos \\
Figuras y cifras de este documento & \tsl{paper/datos3/} de este repositorio \\
Artefactos completos & archivo del autor; se comparten a solicitud (\tsl{contacto@lattimex.com}) \\
\end{tabular}
\end{reproducir}

\bibliographystyle{plain}
\bibliography{referencias}

\end{document}
